% Discusión (consigna 2.2, ítem 9). Base: versión revisada del estudiante (25-nov-2026); los párrafos sobre las pruebas
% econométricas son BORRADOR ASISTIDO POR IA (2026-11-25): reescribir con tus palabras.
\section{Discusión}

El FPF peruano creció con rapidez después de quedar regulado, como predicen \textcite{ran2025} para el crowdfunding de deuda con normas claras. La otra mitad de su resultado, la entrada de nuevas plataformas, no se observa en el Perú: de \cifra{plataformasAutorizadas} sociedades con información en la SMV, solo \cifra{plataformasConFinanciamiento} financiaron proyectos y una concentra casi todo el monto. Quizá influyan la escala y la reputación de las plataformas que actúan como intermediarios \parencite{braggion2025}, algo que con datos agregados no se puede comprobar.

La modificación de noviembre de 2025 coincide con montos mayores. Sin embargo, al descontar la tendencia no aparece un efecto positivo, y varios cortes ficticios anteriores también resultan significativos. Por ello H1 se confirma solo en sentido descriptivo, y con \cifra{mesesDespues} meses posteriores el efecto de la norma tendrá que evaluarse con una serie más larga.

La ausencia de cointegración y de un efecto conjunto de las tasas en el corto plazo indica que, con los 69 meses, no se observa que el monto del FPF se mueva con el precio del crédito tradicional. Estos datos no distinguen entre la sustitución que hallan \textcite{kim2026} y \textcite{bennaceur2026} y la complementariedad que documentan \textcite{beaumont2026} y \textcite{eid2024}. El resultado es sensible a la muestra: con los meses de actividad, Engle-Granger encuentra una relación con la tasa de los bancos a pymes. También pesa que varias tasas tienen quiebres, y que el FPF es un mercado joven cuyo crecimiento puede responder más a la expansión de una plataforma que a las tasas.

En cuanto al acceso, la ventaja de los proyectos empresariales se reduce al controlar la plataforma y el año, y dentro de la líder no es significativa. Casi todos los personales provienen de una plataforma que dejó de registrar financiamiento en 2024. Desde entonces, el FPF peruano es un mercado de préstamos a empresas con montos crecientes, como en los países donde las plataformas amplían el financiamiento sin llegar necesariamente a los prestatarios más vulnerables \parencite{fareed2025}. La Unión Europea tomó otro camino y excluyó a los consumidores como receptores \parencite{morello2025}.

De las plataformas que financiaron proyectos, solo la líder es rentable. La mediana del ROA de las plataformas es menor que la de otros intermediarios supervisados, pero la diferencia no es significativa al controlar el tamaño, en línea con la eficiencia desigual de las fintech de crédito \parencite{onorato2024,pampurini2024}. Con los datos públicos tampoco se puede saber si el FPF llega a las empresas con menos garantías, que según \textcite{hau2024} son las que más se benefician del crédito fintech.

Por la misma razón, la comparación con el crédito tradicional es limitada. Sus tasas son altas y, en varios segmentos, crecientes, pero sin una tasa observada del FPF no se puede afirmar que este sea más barato o más caro. En ese punto la literatura tampoco es concluyente \parencite{chu2024,fareed2025,vives2025}.

Las limitaciones principales son la falta de tasa, plazo, sector y morosidad por operación, el dominio de una sola plataforma, la corta vida del mercado (\cifra{mesesActividadFPF} meses con actividad) y el bajo número de plataformas en la comparación de rentabilidad. Las pruebas son asociativas: un resultado no significativo no equivale a ausencia de relación. Con datos por operación y una serie más larga, futuras investigaciones podrán comparar precios, evaluar la modificación de 2025 e incorporar la demanda, como hacen \textcite{alvarezfranco2026} para América Latina.
